% Partie : théorie de la démonstration
% Chapitre : calcul des séquents
% Section : quantificateurs

\documentclass[../../../include/open-logic-section]{subfiles}

\begin{document}

\olfileid[fr]{pt}{seq}{qua}

\olsection{Preuves régulières et substitution}

Les règles quantifiées présentent quelques difficultés
liées aux conditions sur les constantes propres des
règles \LeftR{\lexists} et~\RightR{\lforall}.
La plupart se résolvent en veillant à ne jamais
réutiliser la constante~$a$ de ces inférences.
Introduisons la définition suivante.

\begin{defn}
Une preuve~$\pi$ de \Log{G1c} ou de \Log{G3c}
est \emph{régulière} si :
\begin{enumerate}
  \item aucune constante~$a$ ne sert de constante propre
  à plus d'une inférence \LeftR{\lexists}
  ou~\RightR{\lforall} ;
  \item si $a$ est la constante propre d'une inférence
  \LeftR{\lexists} ou~\RightR{\lforall},
  ses occurrences dans~$\pi$ se trouvent uniquement
  dans la sous-preuve située au-dessus de cette inférence.
\end{enumerate}
\end{defn}

\begin{lem}\ollabel{lem:subst-regular}
Soit $\pi$ une preuve régulière de séquent final
$\Gamma \Sequent \Delta$.
Soit $c$ une constante qui n'est pas une constante
propre de~$\pi$, et soit $t$ un terme clos ne contenant
aucune constante propre de~$\pi$.
L'arbre $\Subst{\pi}{t}{c}$ obtenu en remplaçant
chaque occurrence de~$c$ par~$t$ est une preuve régulière de même hauteur.
Son séquent final est
$\Subst{\Gamma}{t}{c} \Sequent \Subst{\Delta}{t}{c}$,
où
$\Subst{\Gamma}{t}{c} = \Setabs{!A(t)}{!A(c) \in \Gamma}$.
Cette notation désigne ici un multiensemble :
la substitution s'applique à chaque occurrence
et en conserve la multiplicité. Elle ne fusionne
pas les occurrences devenues identiques.
La même convention vaut pour le conséquent.
\end{lem}

\begin{proof}
  On raisonne par récurrence sur $\pheight{\pi}$.
  Si $\pheight{\pi} = 0$, la preuve se réduit à un axiome.
  La substitution $\Subst{\pi}{t}{c}$ donne encore
  un axiome : toute formule~$!A(c)$ devient $!A(t)$,
  et les deux types d'axiomes conservent leur forme.

  Si $\pheight{\pi} > 0$, on distingue les cas selon
  la dernière inférence de~$\pi$.
  Les cas des règles structurelles, dans \Log{G1c},
  et des règles propositionnelles sont laissés en exercice.
  Considérons les règles quantifiées.
  \begin{enumerate}
    \item Si la dernière inférence est $\RightR{\lforall}$,
    la preuve $\pi$ a la forme
    \[
    \AxiomC{}
    \RightLabel{$\pi'$}
    \Deduce$\Gamma(c) \fCenter \Delta'(c), !A(a,c)$
    \RightLabel{\RightR{\lforall}}
    \UnaryInf$\Gamma(c) \fCenter \Delta'(c), \lforall[x][!A(x,c)]$
    \DisplayProof
    \]
    Par hypothèse de récurrence, $\Subst{\pi'}{t}{c}$
    est une preuve régulière de
    $\Gamma(t) \fCenter \Delta'(t), !A(a,t)$.
    Puisque $a$ est une constante propre de~$\pi$,
    elle ne figure pas dans~$t$.
    La dernière inférence de $\Subst{\pi}{t}{c}$,
    \[
    \AxiomC{}
    \RightLabel{$\Subst{\pi'}{t}{c}$}
    \Deduce$\Gamma(t) \fCenter \Delta'(t), !A(a,t)$
    \RightLabel{\RightR{\lforall}}
    \UnaryInf$\Gamma(t) \fCenter \Delta'(t), \lforall[x][!A(x,t)]$
    \DisplayProof
    \]
    est donc conforme à la règle : sa conclusion
    ne contient pas~$a$.
    \item Si la dernière inférence est $\LeftR{\lforall}$
    dans \Log{G3c}, la preuve $\pi$ a la forme
    \[
    \AxiomC{}
    \RightLabel{$\pi'$}
    \Deduce$!A(s(c),c), \lforall[x][!A(x,c)], \Gamma(c) \fCenter \Delta'(c)$
    \RightLabel{\LeftR{\lforall}}
    \UnaryInf$\lforall[x][!A(x,c)], \Gamma(c) \fCenter \Delta'(c)$
    \DisplayProof
    \]
    Par hypothèse de récurrence, $\Subst{\pi'}{t}{c}$
    est une preuve régulière du séquent
    \[
    !A(s(t),t), \lforall[x][!A(x,t)], \Gamma(t) \fCenter \Delta'(t)\,.
    \]
    Sa dernière inférence,
    dans $\Subst{\pi}{t}{c}$,
    \[
    \AxiomC{}
    \RightLabel{$\Subst{\pi'}{t}{c}$}
    \Deduce$!A(s(t),t), \lforall[x][!A(x,t)], \Gamma(t) \fCenter \Delta'(t)$
    \RightLabel{\LeftR{\lforall}}
    \UnaryInf$\lforall[x][!A(x,t)], \Gamma(t) \fCenter \Delta'(t)$
    \DisplayProof
    \]
    est conforme à la règle. Le cas de \LeftR{\lforall}
    dans~\Log{G1c} est analogue, mais plus simple.
    Le terme d'instanciation reste clos après substitution.
    \item Si la dernière inférence est \RightR{\lexists}
    ou \LeftR{\lexists}, compléter l'argument en exercice.
  \end{enumerate}
  Dans chaque cas, on prolonge une preuve régulière,
  ou deux dans le cas d'une règle à deux prémisses.
  Le nouveau séquent final ne diffère de celui de $\pi$
  que par le remplacement de $c$ par~$t$.
  Les constantes propres de $\pi$ et de
  $\Subst{\pi}{t}{c}$ sont les mêmes :
  aucune n'est remplacée, et le terme substitué
  n'en introduit aucune occurrence nouvelle.
  Chaque constante propre conserve donc son unique
  inférence et toutes ses occurrences restent dans
  la sous-preuve au-dessus de cette inférence.
  Ainsi, $\Subst{\pi}{t}{c}$ est régulière.
  La substitution conserve aussi la structure de l'arbre,
  et donc sa hauteur.
\end{proof}

\begin{prop}
Toute preuve $\pi$ de \Log{G1c} ou de \Log{G3c}
possède une preuve régulière~$\pi'$ de même structure,
et en particulier de même hauteur, qui n'en diffère
que par des renommages de constantes propres.
\end{prop}

\begin{proof}
Pour cette démonstration seulement, appelons \emph{propre}
une inférence à constante propre de~$\pi$
si sa constante ne sert à aucune autre inférence
et n'apparaît nulle part dans~$\pi$ en dehors
de sa sous-preuve immédiate. Appelons-la \emph{sale}
dans le cas contraire.
Soit $n$ le nombre d'inférences sales de~$\pi$.
On raisonne par récurrence forte sur~$n$.
Si $n=0$, toutes les inférences concernées sont
propres et $\pi$ est déjà régulière.
Sinon, choisissons une inférence sale qui n'a
aucune inférence sale au-dessus d'elle.
Supposons qu'il s'agisse de \RightR{\lforall} ;
le cas de \LeftR{\lexists} est analogue.
La sous-preuve~$\pi_1$ qui se termine par
cette inférence a la forme
    \[
    \AxiomC{}
    \RightLabel{$\pi_2$}
    \Deduce$\Gamma \fCenter \Delta, !A(c)$
    \RightLabel{\RightR{\lforall}}
    \UnaryInf$\Gamma \fCenter \Delta, \lforall[x][!A(x)]$
    \DisplayProof
    \]
La condition sur la constante propre garantit
que $c$ n'apparaît ni dans~$\Gamma$,
ni dans~$\Delta$, ni dans $\lforall[x][!A(x)]$.
Toutes les inférences à constante propre situées
dans $\pi_2$ sont propres dans la preuve entière.
La sous-preuve $\pi_2$ est donc régulière,
même si elle contient de telles inférences.
La constante $c$ n'y sert pas de constante propre :
une inférence supérieure utilisant aussi $c$
serait sale, puisqu'elle partagerait sa constante
avec l'inférence choisie.

Choisissons une constante $a$ qui n'apparaît
nulle part dans~$\pi$ et n'est utilisée
par aucune de ses inférences.
Par le \olref{lem:subst-regular}, $\Subst{\pi_2}{a}{c}$
est une preuve régulière de
$\Gamma \fCenter \Delta, !A(a)$.
On peut lui appliquer~$\RightR{\lforall}$ :
la nouvelle constante est absente de la conclusion.
Remplaçons alors $\pi_1$ dans~$\pi$ par $\pi_1'$,
    \[
    \AxiomC{}
    \RightLabel{$\Subst{\pi_2}{a}{c}$}
    \Deduce$\Gamma \fCenter \Delta, !A(a)$
    \RightLabel{\RightR{\lforall}}
    \UnaryInf$\Gamma \fCenter \Delta, \lforall[x][!A(x)]$
    \DisplayProof
    \]
Le séquent final de cette sous-preuve est inchangé,
ainsi que sa structure et sa hauteur.
L'inférence choisie est devenue propre :
sa constante $a$ n'est employée que par elle
et n'apparaît que dans sa sous-preuve supérieure.
Le remplacement de $c$ par~$a$ ne rend sale
aucune inférence auparavant propre, car il n'introduit
aucune occurrence d'une autre constante propre.
Le nombre d'inférences sales a donc strictement
diminué ; il est au plus~$n-1$, et peut être moindre.
L'hypothèse de récurrence s'applique à la preuve
obtenue et donne le résultat.
\end{proof}

Nous ne citerons pas explicitement cette proposition
à chaque emploi. Nous supposerons désormais les
preuves régulières : si nécessaire, un renommage
de leurs constantes propres les met sous cette forme.

\end{document}
